\subsection{Structural definition and role of the Planck constants}
\label{sec:ii-definicion-planck}

\begin{definition}[Structural determination of angular action]
On the oriented line \(\mathcal L_S\), the HMT determination of
angular action is the pair of sections
\(\hbar_{\mathrm{pre}},\hbar_{\mathrm{ret}}\) in
\eqref{eq:el-secciones-accion}, together with their identification by return.
Its structure combines the valuation of the local norm
\(\varphi\mathrm N_{G_H}(\alpha)\), the character \(H_5\), and the
regional continuation \(\mathcal C_\pi\). The action per complete turn
is \(h_a=2\pi\hbar_a\) in each orientation \(a\).
\end{definition}

This definition preserves three successive determinations. The local
quotient in \eqref{eq:el-h4} fixes the sixteen sheets of the norm;
the inequality \eqref{eq:el-decada} fixes its decimal order. The character and
the normalized renewal in \eqref{eq:revision-renovacion} then determine
the pre-return and post-return coordinates. Uniqueness of the tail
uses its first weight and its concatenation rule; the distribution of
the coefficients of \(H_5\) preserves the character declared in
Section~\ref{sec:revision-planck}. On the positive branch, the multiplicative
information of the return is the invariant
\begin{equation}
 \frac{\hbar_{\mathrm{pre}}}{\hbar_{\mathrm{ret}}}
 =R_{\mathrm{act}}=\frac{H_5}{\eta_{\mathrm{ret}}}>1.
 \label{eq:ii-definicion-planck-retorno}
\end{equation}

Its physical role is expressed through the contraction between action and
frequency: a section \(\hbar_a\) assigns energy to an angular
frequency, whereas \(h_a\) expresses the action of one turn. The quotient
\eqref{eq:ii-definicion-planck-retorno} preserves the structural
modification between the two readings even when the common basis
of units is changed. Identification with the physical constant and its precision
are examined in the chart of Section~\ref{sec:revision-planck-contraste},
with the two return coordinates distinguished.
